\section{Introduction}\label{sec:intro}

The pioneers are the ones with arrows in their backs. The claim has been
tested on named markets and naming choices
\citep{GolderTellis1993, SemadeniAnderson2010}, but not on a record where
being early can be read directly from what every entrant wrote, because
that record did not exist in usable form. It does now. Every US trademark
application since 1884 is public, with its goods/services description,
its owner, its counsel of record, and a dated prosecution history that
records whether the mark registered, whether its owner proved five years
later that the product was still in commerce, and whether it was renewed
at ten. This paper re-parses all 13.99 million case files into that dated
form and asks of each filing two questions about its language: how
unusual was this description for its industry, and was it written in
words the industry was moving toward or words it was leaving behind. The
first is the filing's \emph{atypicality}, the second its \emph{lead}. The
paper then follows the filing, and the firm behind it, through the gates
the record provides.

Section~\ref{sec:related} places the exercise: in the trademarks-as-data
literature, which counts filings but has not read their text against
time; in the text-measurement literature the construction descends from;
and in the pioneering and optimal-distinctiveness literatures whose
question it answers at census scale --- including the nomenclature, since
several established terms assert mechanisms this design cannot show.
Section~\ref{sec:measure} builds the measure: a description is reduced to
the themes it draws on, fitted once over the corpus by a topic model; its
theme proportions are compared, by Kullback--Leibler divergence, to the
average proportions of its own Nice class in the five years before its
filing date and the five years after. The average of the two divergences
is atypicality; their difference is lead, written $L$. The section states
why words are not scored directly, why surprise is measured within a
class and never across classes, and the record's gates and counting
conventions.

Sections~\ref{sec:registration}--\ref{sec:success} carry the results.
Registration rises with atypicality and turns down only in the most
atypical decile, is nearly flat in lead --- a quantity no examiner could
observe, since it is computed from filings not yet made --- and
applicants who refile an abandoned application rewrite toward the class's
typical description, so every registered description is in part a
conformed one (Section~\ref{sec:registration}). The central result is the
five-year proof of continued use: among 3.10 million registrations, marks
whose language most led their industry were cancelled $1.4$ points more
often than marks whose language most lagged it, within Nice class and
registration year with drafting and counsel held --- $2.3$ points raw, in
33 of 44 classes, robust to the timing of failure, to who files, to
atypicality strata, to text-duplication, and to three partitions of the
vocabulary; atypical marks survive the proof more often, but that
protection belongs mostly to filings extending a theme established
elsewhere into a new class (Section~\ref{sec:gate}). The penalty is
episodic where language moved fastest: two to three times larger for
marks filed in the late 1990s, reversed for 2000--2004 filings, and
returned from 2008, with the swing inside four technology classes and,
within them, in which themes each era's leading filings were drawn from
(Section~\ref{sec:waves}). Following the owner instead of the mark, the
debut's language does not predict a priced private round; SEC reporting
and listing fall in lead; and the firms that came to hold most of the
linked value described their first product in language their industry
already used (Section~\ref{sec:success}).

Section~\ref{sec:discussion} reads the stages together on one harmonized
sample and states the implications for trademark-based measurement.
Section~\ref{sec:robust} reports what changes under other choices of
resolution, weighting, window, and period; one choice, the decayed
reference window, shifts the headline estimate, by about a seventh. The
full working-paper treatment, an interactive per-class appendix, and
every artifact behind every number are public.

Being early with the language of a product costs at the five-year proof
of continued use, in three of every four product classes and under every
partition of the vocabulary, and it costs most in the industries and the
years in which the language was moving fastest. Being unusual is rewarded
at the same proof, mostly where a theme is imported from elsewhere, and
it is the examiner, not the market, that rewards middling conformity at
registration. The firms that thrived were not the pioneers of language.
They were written in the words their industries already used, and they
are few.
